% !TEX encoding = UTF-8 Unicode
% !TEX TS-program = XeLaTeX
% !TEX root = ParliamentUserGuide.tex

%%%%%%%%%%%%%%%%%%%%%%%%%%%%%%%%%%%%%%%%%%%%%%%%%%
\chapter{Building \acl{pmnt}}
\label{chapter-building-parliament}

\ac{pmnt} is a cross-platform, mixed-language library. It's core is written in portable C++, but it also has a Java interface. As a result of both the cross-platform and multi-language requirements, the build infrastructure for \ac{pmnt} requires a little bit of work to configure. This chapter is your guide through that process.

\acp{pmnt} build infrastructure has two main parts. The top-level portion is based on Gradle, a build tool used in the Java development community. This portion of the infrastructure builds the Java half of the \ac{pmnt} code base, and it also invokes the second portion, which is based on Boost.Build. Boost.Build is a system that is well-adapted to building C++ code. It has the advantages of being portable and much simpler to use than make files. It is also the standard build system of the Boost project, whose libraries are used by the C++ portion of \ac{pmnt}.

This chapter will step through the libraries and tools that \ac{pmnt} depends upon and show you how to configure them on your system. At the end of this chapter, you should have a working copy of the \ac{pmnt} git repository from which you can build the \ac{pmnt} binaries.

%%%%%%%%%%%%%%%%%%%%%%%%%%%%%%%%%%%%%%%%%%%%%%%%%%
\section{Platforms and Prerequisites}

You will need an appropriate C++ compiler for each operating system on which you wish to build. \ac{pmnt} has been tested on the platform and compiler combinations shown in Table~\ref{platforms-and-compilers}. The last column shows the corresponding Boost.Build toolset name, which will appear in the sections below as we configure the \ac{pmnt} build infrastructure.

\begin{table}[htbp]
	\centering
	\begin{tabular}{lll}
		\toprule
		\textbf{Platform} & \textbf{Compiler} & \textbf{Toolset} \\
		\headingrule
		Windows 11 & Visual Studio 2022 (Visual C++ 14.3) & msvc \\
		\midrule
		Mac OS X 26 Tahoe & Xcode 26.5 & clang \\
		\midrule
		Ubuntu 24 & GCC 14.2 & gcc \\
		\midrule
		RHEL 10 & GCC 14.2 & gcc \\
		\bottomrule
	\end{tabular}
	\caption{Supported Platforms and Compilers}
	\label{platforms-and-compilers}
\end{table}

On Macintosh, \ac{pmnt} builds as a universal binary (for both Apple Silicon and Intel chips) using Apple's Xcode development tools.

\ac{pmnt} requires the Java Developer Kit (JDK), version 25:
\begin{itemize}
	\item On Windows and Mac OS, download the JDK installer from Oracle\urlcite{https://www.oracle.com/java/technologies/}.

	\item On Linux, you may need to install one or more packages, depending on your distribution. On Ubuntu, the package you need is ``openjdk-25-jdk-headless.''
\end{itemize}

Finally, you need a client for the git version control system\urlcite{https://git-scm.com/downloads}.

Once you have these prerequisites installed, you can clone the \ac{pmnt} code base from here:
\begin{quote}
	\url{https://github.com/\githubOrg/parliament}
\end{quote}
Because GitHub limits repository size, this does not include larger binary files that are required by the build. To supply these, download this archive:
\begin{quote}\sloppy
	\url{https://github.com/\githubOrg/parliament/releases/download/dependencies-«latest-date»/parliament-dependencies-«latest-date».zip}
\end{quote}
and expand it in the \path|dependencies| subdirectory of your clone. (You can find the latest date by viewing the Releases page for \ac{pmnt} on GitHub.)



%%%%%%%%%%%%%%%%%%%%%%%%%%%%%%%%%%%%%%%%%%%%%%%%%%
\section{Building RocksDB}
\label{sec:BuildingRocksDB}

\ac{pmnt} uses RocksDB\urlcite{https://rocksdb.org/}, an embedded persistent key-value store from Meta. RocksDB is dual-licensed under both the GPLv2 and Apache 2.0 Licenses. The following procedures are based on RocksDB version \rocksdbversion.

\subsection{Building RocksDB on Windows}
\label{sec:BuildingRocksDBOnWindows}

On Windows, we use Microsoft's \path|vcpkg| package manager to download and create the RocksDB library and its dependencies. In your working copy, change to the \path|rocksdb-windows| directory. If you are running Windows on an Intel- or AMD-based machine, issue these commands from a Visual Studio terminal:

\begin{verbatim}
vcpkg install --triplet x64-windows
vcpkg install --triplet x64-windows-static
vcpkg install --triplet x86-windows
vcpkg install --triplet x86-windows-static
\end{verbatim}

This takes a while to execute, so maybe pop a pizza in the oven.

If you are running Windows on an ARM-based machine, issue these commands instead:

\begin{verbatim}
vcpkg install --triplet arm64-windows
vcpkg install --triplet arm64-windows-static
\end{verbatim}

\red{[Something about setting up environment variables or \path|site-config.jam|.]}

Start the Developer Command Prompt window that is a part of Visual Studio installation. Run the build commands from within the Developer Command Prompt window to have paths to the compiler and runtime libraries set.


\subsection{Building RocksDB on Macintosh}
\label{sec:BuildingRocksDBOnMacintosh}

On Mac OS, RocksDB and its dependencies are installable via HomeBrew. However, Parliament wants universal binaries, which we need to build ourselves. This section details alternate instructions to create RocksDB libraries for use with \ac{pmnt}, starting with RocksDB's dependencies. To build RocksDB on Macintosh, first be sure you are using the latest version of Xcode (since RocksDB requires C++20 support) and install CMake\urlcite{https://cmake.org/}.

\subsubsection{Building Snappy}
\label{sec:BuildingSnappy}

Snappy\urlcite{https://github.com/google/snappy} is an optional dependency of RocksDB. Because its build infrastructure cannot support universal binaries, we omit it.

\subsubsection{Building GFlags}
\label{sec:BuildingGFlags}

Download GFlags\urlcite{https://gflags.github.io/gflags/} version 2.3.0 and then execute these commands:

{\small\begin{verbatim}
tar xf gflags-2.3.0.tar.gz
cd gflags-2.3.0
cmake -S . -B build-cmake -DCMAKE_INSTALL_PREFIX=/opt/gflags-2.3.0 \
   -DBUILD_SHARED_LIBS=ON -DBUILD_STATIC_LIBS=ON \
   -DCMAKE_OSX_ARCHITECTURES="x86_64;arm64"
cd build-cmake
make
sudo make install
\end{verbatim}}

\subsubsection{Building LZ4}
\label{sec:BuildingLZ4}

Download LZ4\urlcite{https://github.com/lz4/lz4} version 1.10.0 and then execute these commands:

{\small\begin{verbatim}
tar xf lz4-1.10.0.tar.gz
cd lz4-1.10.0
make PREFIX=/opt/lz4-1.10.0 CFLAGS="-arch x86_64 -arch arm64" \
   LDFLAGS="-arch x86_64 -arch arm64 -dynamiclib"
sudo make install PREFIX=/opt/lz4-1.10.0
\end{verbatim}}

\subsubsection{Building ZStd}
\label{sec:BuildingZStd}

Download ZStd\urlcite{https://github.com/facebook/zstd} version 1.5.7 and then execute these commands:

{\small\begin{verbatim}
untar xf zstd-1.5.7.tar.gz
cd zstd-1.5.7
make PREFIX=/opt/zstd-1.5.7 CFLAGS="-arch x86_64 -arch arm64" \
   LDFLAGS="-arch x86_64 -arch arm64 -dynamiclib"
sudo make install PREFIX=/opt/zstd-1.5.7
\end{verbatim}}

\subsubsection{Building RocksDB Itself}
\label{sec:BuildingRocksDBItself}

Define the following environment variables so that the builds can find the headers and libraries:
{\small\begin{verbatim}
export GFLAGS_HOME=/opt/gflags-2.3.0
export LZ4_HOME=/opt/lz4-1.10.0
export ROCKSDB_HOME=/opt/rocksdb-11.1.1
export ZSTD_HOME=/opt/zstd-1.5.7
\end{verbatim}}

Now download and expand the source code archive for RocksDB version \rocksdbversion{} from the releases page of the RocksDB GitHub site\urlcite{https://github.com/facebook/rocksdb/releases}. Change line 1541 of util/compression.cc from this:

{\small\begin{verbatim}
Status ExtractUncompressedSize(Args& args) override {
\end{verbatim}}

to this:

{\small\begin{verbatim}
Status ExtractUncompressedSize([[maybe_unused]] Args& args) \
   override {
\end{verbatim}}

On lines 292, 307, 317, 327, 340, 353, 366, and 381 of the file \path|build_tools/build_detect_platform|, replace ``-x'' with ``-c -x''. In addition, on line 381 replace ``-o /dev/null'' with ``-o test.o''.

Comment out lines 355, 356, 360, 363, and 364 of \path|Makefile| by prefixing them with \texttt{\#}.

Execute these commands:

{\small\begin{verbatim}
make static_lib shared_lib PREFIX=/opt/rocksdb-11.1.1 \
   SKIP_DEPENDS=1 EXTRA_CXXFLAGS="-arch x86_64 -arch arm64 \
   -I${GFLAGS_HOME}/include -I${LZ4_HOME}/include \
   -I${ZSTD_HOME}/include" EXTRA_LDFLAGS="-arch x86_64 -arch arm64 \
   -dynamiclib -L${GFLAGS_HOME}/lib -L${LZ4_HOME}/lib \
   -L${ZSTD_HOME}/lib"
sudo make install PREFIX=/opt/rocksdb-11.1.1 SKIP_DEPENDS=1
\end{verbatim}}

Finally, run this test on the resulting dynamic library to be sure the build worked correctly:

{\small\begin{verbatim}
otool -L $ROCKSDB_HOME/lib/librocksdb.11.1.1.dylib
\end{verbatim}}

This should result in output like the following:

{\small\begin{verbatim}
.../librocksdb.11.1.1.dylib (architecture x86_64):
   librocksdb.10.10.dylib ...
   @rpath/libgflags.2.3.dylib ...
   /usr/lib/libz.1.dylib ...
   /usr/lib/libbz2.1.0.dylib ...
   /opt/lz4-1.10.0/lib/liblz4.1.dylib ...
   /opt/zstd-1.5.7/lib/libzstd.1.dylib ...
   /usr/lib/libc++.1.dylib ...
   /usr/lib/libSystem.B.dylib ...
.../librocksdb.11.1.1.dylib (architecture arm64):
   librocksdb.10.10.dylib ...
   @rpath/libgflags.2.3.dylib ...
   /usr/lib/libz.1.dylib ...
   /usr/lib/libbz2.1.0.dylib ...
   /opt/lz4-1.10.0/lib/liblz4.1.dylib ...
   /opt/zstd-1.5.7/lib/libzstd.1.dylib ...
   /usr/lib/libc++.1.dylib ...
   /usr/lib/libSystem.B.dylib ...
\end{verbatim}}

The main points to check for are (a) that the dynamic library contains two architectures (i.e., it is a universal binary) and (b) that all of the optional dependencies (GFlags, LibZ, LibBZ2, LZ4, and ZStd) are linked in.


\subsection{Building RocksDB on Linux}
\label{sec:BuildingRocksDBOnLinux}

To build RocksDB on Linux, upgrade your gcc compiler to at least version 11 to get C++20 support. Next, use the package manager to install RocksDB's dependencies. On Ubuntu, this looks like so:
{\small\begin{verbatim}
sudo apt install libgflags-dev libsnappy-dev zlib1g-dev libbz2-dev \
   liblz4-dev libzstd-dev libxxhash-dev
\end{verbatim}}

Now execute the following commands from the root directory of the expanded RocksDB release:

{\small\begin{verbatim}
make static_lib shared_lib PREFIX=/opt/rocksdb-11.1.1 \
   SKIP_DEPENDS=1
sudo make install PREFIX=/opt/rocksdb-11.1.1 SKIP_DEPENDS=1
\end{verbatim}}

Define the following environment variable so that the \ac{pmnt} build can find the headers and libraries:
{\small\begin{verbatim}
export ROCKSDB_HOME=/opt/rocksdb-11.1.1
\end{verbatim}}

Finally, run this test on the resulting dynamic library to be sure the build worked correctly:

{\small\begin{verbatim}
ldd $ROCKSDB_HOME/lib/librocksdb.so.11.1.1
\end{verbatim}}

This should result in output like the following:

{\small\begin{verbatim}
   /lib/ld-linux-aarch64.so.1 ...
   libbz2.so.1.0 ...
   libc.so.6 ...
   libgcc_s.so.1 ...
   libgflags.so.2.2 ...
   liblz4.so.1 ...
   libm.so.6 ...
   libsnappy.so.1 ...
   libstdc++.so.6 ...
   libxxhash.so.0 ...
   libz.so.1 ...
   libzstd.so.1 ...
   linux-vdso.so.1 ...
\end{verbatim}}

The main point to check is that all of the optional dependencies (GFlags, LibZ, LibBZ2, LZ4, and ZStd) are linked in.


%%%%%%%%%%%%%%%%%%%%%%%%%%%%%%%%%%%%%%%%%%%%%%%%%%
\section{Building the Boost Libraries}
\label{sec:BuildingBoost}

Download the Boost source distribution (version \boostversion{.} or later) from the Boost web site. Unpack this to a handy location on your disk and define this environment variable pointing to that location:

{
	%\renewcommand{\arraystretch}{1.5}
	\renewcommand{\tabcolsep}{0pt}
	\begin{tabular}{l@{\hspace{2em}}l}
		\texttt{export BOOST\_ROOT=\char`~/boost\_\boostversion{\_}}
			& (on Mac OS and Linux)\\
		\texttt{set BOOST\_ROOT=C:\textbackslash{}boost\_\boostversion{\_}}
			& (on Windows)\\
	\end{tabular}
}

From \acp{pmnt} point of view, there are two primary components contained within \verb|BOOST_ROOT|. The first (and most obvious) is the boost libraries themselves. Most of these are so-called ``header-only'' libraries, meaning that there is no pre-compiled library code. All of the code of such libraries is referenced via \verb|#include| directives and compiled along with the calling code. Such libraries are extremely convenient, because they require little setup. \ac{pmnt} also uses several Boost libraries that are not header-only. We will discuss how to build these libraries below.

The second major Boost component is Boost.Build. This is a cross-platform build system (located in the \path|BOOST_ROOT/tools/build|) that is written in a specialized interpreted language whose interpreter is a command called \verb|b2|. The Boost community does not provide \verb|b2| binaries. Rather, the Boost distribution contains a bootstrapping script that builds \verb|b2| from source on your platform.

To build the minimal set of libraries required for \ac{pmnt}, follow the directions below for your platform. In each case, you will invoke the bootstrap script to create the Boost.Build interpreter, and then you will invoke that to build the libraries.

\subsection{Building Boost on Windows}

First, identify (or create) a directory where we will place the \path|b2| build utility and define these environment variables:

{\small\begin{verbatim}
set B2_HOME=«your-directory-here»
set PATH=%PATH%;%B2_HOME%\b2
\end{verbatim}}

Open a Command Prompt, change to the \verb|BOOST_ROOT| directory, and issue these commands to build \path|b2|:

{\small\begin{verbatim}
pushd tools\build
rd /s/q "%B2_HOME%\b2"
.\bootstrap.bat msvc & .\b2 install --prefix="%B2_HOME%\b2"
popd
\end{verbatim}}

On Intel-based computers, issue this command to build the boost libraries:

{\small\begin{verbatim}
b2 -q --disable-icu --ignore-site-config --layout=versioned
--build-dir=build-msvc --stagedir=stage-msvc --with-atomic
--with-chrono --with-container --with-date_time --with-filesystem
--with-log --with-regex --with-serialization --with-test
--with-thread define=BOOST_LOG_USE_STD_REGEX
define=BOOST_LOG_WITHOUT_EVENT_LOG define=BOOST_LOG_WITHOUT_IPC
define=BOOST_LOG_WITHOUT_ASIO define=BOOST_LOG_WITHOUT_DEBUG_OUTPUT
define=BOOST_LOG_WITHOUT_SYSLOG
define=BOOST_TEST_ALTERNATIVE_INIT_API toolset=msvc
address-model=64 variant=debug,release link=shared,static
runtime-link=shared architecture=x86 cxxstd=20 stage
\end{verbatim}}

If you are using Windows on an ARM-based computer, use the same command, but change \path|architecture=x86| to \path|architecture=arm|.

This will build the libraries in \path|stage-msvc/lib|. The following commands will delete intermediate build products to save disk space:

{\small\begin{verbatim}
rd /s/q build-msvc
del tools\build\b2.exe tools\build\src\engine\b2.exe
\end{verbatim}}

\subsection{Building Boost on Macintosh}

First, set these environment variables:

{\small\begin{verbatim}
export B2_HOME=/opt/b2
export PATH=$B2_HOME/bin:$PATH
\end{verbatim}}

Open a shell, change to the \verb|BOOST_ROOT| directory, and issue these commands to build \path|b2|:

{\small\begin{verbatim}
pushd tools\build
./bootstrap.sh ; ./b2 install --prefix=temp-boost-build
sudo rm -r /opt/b2 ; sudo cp -R temp-boost-build /opt/b2
popd
\end{verbatim}}

Now issue this command to build the boost libraries:

{\small\begin{verbatim}
b2 -q --disable-icu --ignore-site-config --layout=versioned
--build-dir=build-clang --stagedir=stage-clang --with-atomic
--with-charconv --with-chrono --with-container --with-date_time
--with-filesystem --with-log --with-regex --with-serialization
--with-test --with-thread define=BOOST_LOG_USE_STD_REGEX
define=BOOST_LOG_WITHOUT_SYSLOG define=BOOST_LOG_WITHOUT_IPC
define=BOOST_LOG_WITHOUT_ASIO
define=BOOST_TEST_ALTERNATIVE_INIT_API toolset=clang
address-model=64 variant=debug,release link=shared,static
runtime-link=shared architecture=arm+x86 cxxstd=20 stage
\end{verbatim}}

This will build the libraries in \path|stage-clang/lib|. The following command will delete intermediate build products to save disk space:

{\small\begin{verbatim}
rm -r tools/build/b2 tools/build/src/engine/b2
tools/build/temp-boost-build build-clang
\end{verbatim}}

\subsection{Building Boost on Linux}
\label{sec:BuildingBoostOnLinux}

First, set these environment variables:

{\small\begin{verbatim}
export B2_HOME=/opt/b2
export PATH=$B2_HOME/bin:$PATH
export LINUX_DISTRO=«platform»
\end{verbatim}}

where «platform» indicates the particular flavor of Linux you are using. For the flavors of Linux that are officially supported, this should be set as given in Table~\ref{tbl:SupportedLinuxFlavors}.

\begin{table}[htbp]
	\centering
	\begin{tabular}{ll}
		\toprule
		\emph{Platform}	& \emph{Description}\\
		\headingrule
		\path|rhel10|		& Red Hat Enterprise Linux 10\\
		\path|ubuntu26|	& Ubuntu Linux 26, LTS\\
		\path|ubuntu24|	& Ubuntu Linux 24, LTS\\
		\bottomrule
	\end{tabular}
	\caption{Supported Linux Flavors}
	\label{tbl:SupportedLinuxFlavors}
\end{table}

Open a shell, change to the \verb|BOOST_ROOT| directory, and issue these commands to build \path|b2|:

{\small\begin{verbatim}
pushd tools\build
./bootstrap.sh ; ./b2 install --prefix=temp-boost-build
sudo rm -r /opt/b2 ; sudo cp -R temp-boost-build /opt/b2
popd
\end{verbatim}}

On Intel-based computers, issue this command to build the boost libraries:

{\small\begin{verbatim}
b2 -q --disable-icu --ignore-site-config --layout=versioned
--build-dir=build-msvc --stagedir=stage-msvc --with-atomic
--with-chrono --with-container --with-date_time --with-filesystem
--with-log --with-regex --with-serialization --with-test
--with-thread define=BOOST_LOG_USE_STD_REGEX
define=BOOST_LOG_WITHOUT_EVENT_LOG define=BOOST_LOG_WITHOUT_IPC
define=BOOST_LOG_WITHOUT_ASIO define=BOOST_LOG_WITHOUT_DEBUG_OUTPUT
define=BOOST_LOG_WITHOUT_SYSLOG
define=BOOST_TEST_ALTERNATIVE_INIT_API toolset=msvc
address-model=64 variant=debug,release link=shared,static
runtime-link=shared architecture=x86 cxxstd=20 stage
\end{verbatim}}

If you are using Linux on an ARM-based computer, use the same command, but change \path|architecture=x86| to \path|architecture=arm|.

This will build the libraries in \path|stage-gcc-$LINUX_DISTRO/lib|. Run these commands to delete intermediate build products and save disk space:

{\small\begin{verbatim}
rm -r tools/build/b2 tools/build/src/engine/b2
tools/build/temp-boost-build build-gcc-$LINUX_DISTRO
\end{verbatim}}


%%%%%%%%%%%%%%%%%%%%%%%%%%%%%%%%%%%%%%%%%%%%%%%%%%
\section{Configuring Boost.Build}

Next we need to configure Boost.Build so that it can build \ac{pmnt}. Start by defining the following environment variables:
\begin{itemize}
	\item\verb|LINUX_DISTRO|: \emph{(Linux only)} The flavor of Linux you are using. See Table~\ref{tbl:SupportedLinuxFlavors} in Section~\ref{sec:BuildingBoostOnLinux} above.

	\item\verb|JAVA_HOME|: The location of your JDK installation, which must be version 25 or higher.

	\item\verb|BOOST_VERSION|: The version of Boost (currently \boostversion{\_}).

	\item\verb|BOOST_ROOT|: As described above in Section~\ref{sec:BuildingBoost}.

	\item\verb|BOOST_BUILD_PATH|: The sub-directory \path|tools/build| of \verb|BOOST_ROOT|.

	\item\verb|BOOST_TEST_LOG_LEVEL|: Controls the output volume of the C++ unit tests. Possible values and their meanings are listed in Table~\ref{boost-test-log-level-values}.

	\begin{table}[htbp]
		\centering
		\begin{tabular}{ll}
			\toprule
			\textbf{Value} & \textbf{Meaning} \\
			\headingrule
			\verb|all|				& report all log messages \\
			\verb|success|			& the same as all \\
			\verb|test_suite|		& show test suite messages \\
			\verb|message|			& show user messages \emph{(useful default)} \\
			\verb|warning|			& report warnings issued by user \\
			\verb|error|			& report all error conditions \\
			\verb|cpp_exception|	& report uncaught C++ exceptions \\
			\verb|system_error|	& report system-originated non-fatal errors \\
			\verb|fatal_error|	& report only fatal errors \\
			\verb|nothing|			& do not report any information \\
			\bottomrule
		\end{tabular}
		\caption{Possible Values of \texttt{BOOST\_TEST\_LOG\_LEVEL}}
		\label{boost-test-log-level-values}
	\end{table}
\end{itemize}

On Windows only, define the following environment variables as well:

{\footnotesize\begin{verbatim}
set VS_ROOT=C:\Program Files\Microsoft Visual Studio\2022\Community\VC\
set VS_BAT_PATH_X86_64=!VS_ROOT!Auxiliary\Build\vcvars64.bat
set VS_BAT_PATH_ARM_64=!VS_ROOT!Auxiliary\Build\vcvarsarm64.bat
set VS_REDIST_PATH_X86_64=!VS_ROOT!Redist\MSVC\v143\vc_redist.x64.exe
set VS_REDIST_PATH_ARM=!VS_ROOT!Redist\MSVC\v143\vc_redist.arm64.exe
\end{verbatim}}

Note that on Windows, in a command prompt, you will want to run one of the \path|VS_BAT_PATH_«arch»_64| batch files to set up your command prompt for running the C++ compiler.

Next we create two Boost.Build configuration files, \path|site-config.jam| and \path|user-config.jam|. Boost.Build reads these files on startup. These files can be placed in a number of locations. Table~\ref{boost-build-config-search} explains where Boost.Build searches to find these files.

\begin{table}[htbp]
	\centering
	\begin{tabular}{>{\small}l>{\RaggedRight\small}p{49mm}>{\RaggedRight\small}p{49mm}}
		\toprule
		OS & \path|site-config.jam| & \path|user-config.jam| \\
		\headingrule
		Unix-like:
			&	\path|/etc| \newline
				\verb|$HOME| \newline
				\verb|$BOOST_BUILD_PATH|
			&	\verb|$HOME| \newline
				\verb|$BOOST_BUILD_PATH| \\
		\midrule
		Windows:
			&	\verb|%SystemRoot%| \newline
				\verb|%HOMEDRIVE%%HOMEPATH%| \newline
				\verb|%HOME%| \newline
				\verb|%BOOST_BUILD_PATH%|
			&	\verb|%HOMEDRIVE%%HOMEPATH%| \newline
				\verb|%HOME%| \newline
				\verb|%BOOST_BUILD_PATH%| \\
		\bottomrule
	\end{tabular}
	\caption{Boost.Build Search Paths for Configuration Files}
	\label{boost-build-config-search}
\end{table}

Some people prefer to keep these files together with their Boost.Build installation, in \verb|BOOST_BUILD_PATH|. However, there are example files in that directory already, so you will have to rename them if you choose this option. Others prefer to separate the configuration files from the Boost.Build installation so that they can update to a new version of Boost.Build without having to first save \path|site-config.jam| and \path|user-config.jam| and then restore them after the update is complete. On Windows, using the \verb|HOME| location requires defining the environment variable \verb|HOME|, because Windows does not define it by default.

Within your working copy of the \ac{pmnt} repository, in the directories \path|doc/MacOS|, \path|doc/Windows|, and \path|doc/Linux|, you will find example configuration files for Macintosh, Windows, and Linux respectively that you can copy and customize. The most important customization that you need to make is to remove (or comment out) any lines in \path|user-config.jam| for compiler versions that you have not installed.



%%%%%%%%%%%%%%%%%%%%%%%%%%%%%%%%%%%%%%%%%%%%%%%%%%
\section{Building \ac{pmnt} Itself}

The file \path|build.properties|, located in the root of your working copy, controls the \ac{pmnt} build. This file does not exist by default. If the build does not find it, it uses \path|build.properties.default| instead. The latter contains the build options used to create an official release of \ac{pmnt}, which may include several different release builds. To build a single variant, first copy \path|build.properties.default| to \path|build.properties|. Then customize the latter. The file itself contains instructions. Please change \path|build.properties.default| directly \emph{only if you intend to update the official release build options.}

\red{\emph{This next bit about the top-level build doesn't work just yet, but should be working when we get to the next release.}}

You are now ready to build \ac{pmnt}. From the root directory of your \ac{pmnt} working copy, issue this command:

{
	%\renewcommand{\arraystretch}{1.5}
	\renewcommand{\tabcolsep}{0pt}
	\begin{tabular}{l@{\hspace{2em}}l}
		\path|./gradlew build|		& (on Mac OS and Linux)\\
		\path|.\gradlew.bat build|	& (on Windows)\\
	\end{tabular}
}

This builds the entire repository, including both native and Java code, and creates a distribution-ready package in this directory:
\begin{quote}
	\texttt{target/parliament-\textbf{\textit{x.y.z}}-\textbf{\textit{platform}}}
\end{quote}
where \texttt{\textbf{\textit{x.y.z}}} is the \ac{pmnt} version number and \texttt{\textbf{\textit{platform}}} indicates the platform on which the distribution runs, as shown in Table~\ref{tbl:SupportedPlatforms}.

To delete all build products, issue this command:

{
	%\renewcommand{\arraystretch}{1.5}
	\renewcommand{\tabcolsep}{0pt}
	\begin{tabular}{l@{\hspace{2em}}l}
		\path|./gradlew clean|		& (on Mac OS and Linux)\\
		\path|.\gradlew.bat clean|	& (on Windows)\\
	\end{tabular}
}

%The \path|build.properties| file also contains an option that disables the native code unit tests, \verb|skipNativeUnitTest|. This is important when cross-compiling, e.g., building 64-bit binaries on 32-bit Windows. If you use this option, be sure to change it only in \path|build.properties|, not \path|build.properties.default|.

For more targeted builds, gradle can be run from the sub-projects in your working copy. These include (in the order they should be built):
\begin{itemize}
	\item\path|Parliament|: The native code at the heart of \ac{pmnt}, plus its \ac{jni} interface
	\item\path|jena/ParliamentClient|: A \ac{pmnt}-specific client-side Java library for communicating with a \ac{pmnt} \ac{sparql} endpoint
	\item\path|jena/SparqlQueryBuilder|: A client-side Java library to both issue queries and to dynamically build and modify queries at run time. This library is  \ac{pmnt}-independent.
	\item\path|jena/OntologyDrivenDataAccess|: A client-side Java library that issues queries and returns the results in the form of Java data objects. This library is  \ac{pmnt}-independent.
	\item\path|jena/OntologyBundle|: A Gradle plugin to package (bundle) an \ac{owl} ontology in a jar in order to make it easy to use in Java-based software projects.
	\item\path|jena/KbGraph|: Enables Jena to use \ac{pmnt} for storing models
	\item\path|jena/NumericIndex|: An add-on to KbGraph for processing numeric queries efficiently
	\item\path|jena/TemporalIndex|: An add-on to KbGraph for processing temporal queries efficiently
	\item\path|jena/SpatialIndex|: An add-on to KbGraph for processing spatial queries efficiently
	\item\path|server|: A Spring-based \ac{sparql} endpoint implemented on top of KbGraph.
\end{itemize}




When working on the native code portions of \ac{pmnt}, it can be useful to run the \path|b2| portion of the build directly. To do so, change directory to \path|KbCore| (for the \ac{pmnt} \ac{dll} itself), \path|AdminClient| (for the command line interface to \ac{pmnt}), or \path|Test| (for the unit tests). These directories are located within the \path|Parliament| sub-directory of your working copy. Then issue the command
\begin{verbatim}
b2 -q «build-options»
\end{verbatim}
Here \verb|«build-options»| is a placeholder for one set of build options from \path|build.properties|, described above. The \verb|-q| option causes \path|b2| to quit immediately whenever an error occurs.
